% Part: turing-machines 
% Chapter: undecidability 
% Section: unsolvability-decision-problem

\documentclass[../../../include/open-logic-section]{subfiles}

\begin{document}

\olfileid{tur}{und}{uns} 
\olsection{فیصلے دا مسئلہ حل نہیں ہو سکدا}

\begin{thm}
\ollabel{thm:decision-prob}
فیصلے دا مسئلہ حل نہیں ہو سکدا۔ کوئی ایہو جیہی ٹیورنگ
مشین~$D$ نہیں کہ جدوں پٹی اُتے پہلے درجے دی منطق
دا اک !!a{sentence}~$!B$ انپٹ وجوں رکھ کے اوہنوں
چلایا جاوے، تاں~$D$ آخرکار رُکے تے~$1$ آؤٹ پٹ اُدوں تے صرف
اُدوں دے جدوں~$!B$ درست ہووے، تے دوجی صورت وچ~$0$ دے۔
\end{thm}

\begin{proof}
فرض کرو فیصلے دا مسئلہ حل ہو سکدا ہووے، یعنی
ایہو جیہی ٹیورنگ مشین~$D$ موجود ہووے۔ فیر اسی
رُکّن دا مسئلہ ایویں حل کر سکدے آں۔
اسی اک ٹیورنگ مشین~$E$ بناندے آں جیہڑی
ٹیورنگ مشین دا اشاریہ~$e$، یعنی~$M_e$ دا نمبر، تے انپٹ~$w$
ملن اُتے متعلقہ !!{sentence}~
$!T(M_e, w) \lif !E(M_e, w)$ حساب کردی
اے تے پٹی دا سب توں کھبّا خانہ پڑھدی ہوئی
رُک جاندی اے۔ فیر مشین~$E \concat D$،
انپٹ~$e$ تے~$w$ ملن اُتے، پہلے~
$!T(M_e, w) \lif
!E(M_e, w)$ حساب کرے گی
تے فیر اوس انپٹ اُتے فیصلے دے مسئلے والی
مشین~$D$ چلائے گی۔ $D$ آؤٹ پٹ~$1$ نال
اُدوں تے صرف اُدوں رُکدی اے جدوں
$!T(M_e, w) \lif !E(M_e, w)$ درست ہووے،
ورنہ~$0$ دیندی اے۔
\olref[ver]{lem:halt-if-valid} تے
\olref[ver]{lem:valid-if-halt} مطابق
$!T(M_e, w) \lif !E(M_e, w)$ اُدوں تے صرف اُدوں
درست اے جدوں~$M_e$ انپٹ~$w$ اُتے رُکدی اے۔
ایس لئی~$E\concat D$، انپٹ~$e$ تے~$w$
ملن اُتے، آؤٹ پٹ~$1$ نال اُدوں تے صرف اُدوں رُکدی
اے جدوں~$M_e$ انپٹ~$w$ اُتے رُکدی اے،
تے دوجی صورت وچ آؤٹ پٹ~$0$ نال رُکدی اے۔ دوجے لفظاں
وچ~$E \concat D$ رُکّن دا مسئلہ حل
کر دیوے گی۔ پر~\olref[hal]{thm:halting-problem}
توں سانوں پتا اے کہ ایہو جیہی کوئی ٹیورنگ مشین
موجود نہیں ہو سکدی۔
\end{proof}

\begin{cor}\ollabel{cor:undecidable-sat}%
پہلے درجے دی منطق دا دتا ہویا !!{sentence}
پورا ہو سکدا اے یا نہیں، ایہدا فیصلہ نہیں کیتا جا سکدا۔
\end{cor}

\begin{proof}
  فرض کرو اک ٹیورنگ مشین~$S$ ایہ فیصلہ کر
  سکدی ہووے کہ جملہ پورا ہو سکدا اے یا نہیں۔
  فیر اسی فیصلے دا مسئلہ ایویں حل کر سکدے
  آں۔ اک !!a{sentence}~$B$ انپٹ وجوں ملے تاں
  $!B$ نوں اک خانہ سجّے ول سرکا دیو۔
  خانہ~$1$ اُتے واپس آ کے نشان~$\lnot$ لکھ دیو۔

  ہن ٹیورنگ مشین~$S$ چلاؤ۔ ایہ آخرکار پٹی
  اُتے یا آؤٹ پٹ~$1$ نال رُکے گی (جے
  $\lnot !B$ پورا ہو سکدا ہووے) یا~$0$
  نال (جے~$\lnot !B$ پورا نہ ہو سکدا ہووے)۔
  جے~$\TMstroke$ خانہ~$1$ اُتے ہووے تاں
  اوہ مٹا دیو؛ جے خانہ~$1$ خالی ہووے تاں
  اوتھے~$\TMstroke$ لکھ دیو، فیر رُک جاؤ۔

  ایہ ٹیورنگ مشین ہمیشہ رُکدی اے، تے
  آؤٹ پٹ~$1$ اُدوں تے صرف اُدوں دیندی اے جدوں
  $\lnot
  !B$ پورا نہ ہو سکدا ہووے، ورنہ~$0$۔
  کیوں جے~$!B$ اُدوں تے صرف اُدوں درست اے جدوں
  $\lnot !B$~پورا نہ ہو سکدا ہووے،
  ایس لئی مشین~$1$ اُدوں تے صرف اُدوں دیندی اے
  جدوں~$!B$ درست ہووے، ورنہ~$0$؛ یعنی
  ایہ فیصلے دا مسئلہ حل کر دیوے گی۔
\end{proof}

\begin{explain}
ایس لئی کوئی ایہو جیہی ٹیورنگ مشین نہیں جیہڑی
سوال «کی~$!B$ پہلے درجے دی منطق دا
درست !!{sentence} اے؟» دا ہمیشہ ٹھیک
«ہاں» یا «نہیں» جواب دیوے۔ پر اک ایہو جیہی
ٹیورنگ مشین \emph{موجود اے} جیہڑی جدوں صحیح
جواب «ہاں» ہووے تاں «ہاں» آکھدی اے، تے جدوں
صحیح جواب «نہیں» ہووے تاں رُکدی نہیں۔
ایہ پہلے درجے دی منطق دے صحت مندی تے
تمامیت دے قضیے، تے ایس حقیقت توں نکلدا اے
کہ سارے !!{derivation}s مؤثر طریقے نال
اک اک کر کے فہرست وچ لکھے جا سکدے نیں۔
\end{explain}

\begin{thm}
  \ollabel{thm:valid-ce}%
  پہلے درجے دے !!{sentence}s درست نیں یا نہیں، ایہ
  طے کرن دا مسئلہ نیم فیصلہ کرن جوگ اے۔ اک ٹیورنگ مشین~$E$
  موجود اے جیہڑی پٹی اُتے پہلے درجے دی
  منطق دا اک !!a{sentence}~$!B$ انپٹ وجوں
  رکھ کے شروع کیتی جائے تاں~$E$ آخرکار رُک
  کے آؤٹ پٹ~$1$ اُدوں تے صرف اُدوں دیندی اے
  جدوں~$!B$ درست ہووے۔ دوجی صورت وچ رُکدی نہیں۔
\end{thm}

\begin{proof}
  پہلے درجے دی منطق دے سارے ممکن
  !!{derivation}s اک مؤثر طریقے نال اک
  توں بعد اک بنائے جا سکدے نیں۔ مشین~$E$
  ایہی کردی اے، تے جدوں ایہ اک ایہو جیہا
  !!a{derivation} لبھدی اے جیہڑا
  $\Proves
  !B$ دکھاندا اے، تاں آؤٹ پٹ~$1$ نال
  رُک جاندی اے۔ صحت مندی دے قضیے مطابق
  جے~$E$ آؤٹ پٹ~$1$ نال رُکدی اے تاں
  ایس دی وجہ~$\Entails !B$ اے۔
  تمامیت دے قضیے مطابق جے~$\Entails !B$
  ہووے تاں اک !!a{derivation} موجود اے
  جیہڑا~$\Proves !B$ دکھاندا اے۔ چونکہ
  $E$ سارے ممکن !!{derivation}s ترتیب
  نال بناندی اے، ایس لئی اوہ آخرکار
  $\Proves !B$ دکھان والا اخذ لبھ لوے گی،
  تے آؤٹ پٹ~$1$ نال رُک جائے گی۔
\end{proof}

\end{document}
